\documentclass{article}

\usepackage[utf8x]{inputenc}
\usepackage[T1]{fontenc}
\usepackage[francais]{babel}
\usepackage{lmodern}
\usepackage{stmaryrd}
\usepackage{amsmath, amsthm}
\usepackage{amsfonts, amssymb}
\usepackage{graphicx}
\usepackage[nottoc, notlof, notlot]{tocbibind}
\usepackage{pst-circ, pst-plot, pstricks-add}
\usepackage{array}
\usepackage{url}
\usepackage[colorlinks,linkcolor=black]{hyperref}
\usepackage{ifthen}

\newcommand{\ds}{\displaystyle}

\renewcommand{\ge}{\geqslant}
\renewcommand{\le}{\leqslant}
\renewcommand{\preceq}{\preccurlyeq}
\renewcommand{\succeq}{\succcurlyeq}

\newcommand{\Numero}{\No}
\newcommand{\numero}{\no}

\newcommand{\fixme}{\textbf{FIXME}}

\newcommand{\forplayer}{n}
\newcommand{\theplayer}{1}

\newcommand{\player}[2]{\ifthenelse{\equal{\forplayer}{y}}{\ifthenelse{\equal{\theplayer}{#1}}{#2}{}}{#2}}
\newcommand{\mj}[1]{\ifthenelse{\equal{\forplayer}{n}}{#1}{}}

\title{Les Conspirations d’Anacréon}
\author{Martin \textsc{Bodin}}
\date{}

\begin{document}

\maketitle

\tableofcontents

\newpage

\newlength{\annee}
\settowidth{\annee}{Années 10 000}
\newlength{\texte}
\setlength{\texte}{\textwidth} \addtolength{\texte}{-\annee} 
	\addtolength{\texte}{-2\tabcolsep}

\section{Présentation générale de l’Univers}

Je présente ici très rapidement l’univers, en faisant plusieurs simplifications.
J’espère que personne ne m’en voudra trop…
Si vous voulez des détails sur un point de l’univers, demandez moi !

Cela va faire plusieurs milliers d’années que l’homme a conquis la galaxie.
Il a fait énormément d’avancées technologiques, même si les connaissances tendent à devenir trop complexe et que les avancées sont de plus en plus rares en matière technologique.

Je met ici une frise chronologique de l’humanité : \\
\begin{tabular}{@{}p{\annee}p{\texte}@{}}
Années 2100 & L’homme est capable de se déplacer de manière acceptablement rapide dans tout le système solaire, les premiers robots (métalliques, mais de forme humaines) sont créés. \\
Années 2600 & La Terre combat ses colonies qui réclament leur indépendance. Les colonies gagnent et commencent à former un tout politique. \\
Années 4700 & Les villes terrestres s'enferment sous de gigantesques dômes. Il n’y a plus que des mégalopoles où les hommes ont peur de voir la lumière du Soleil. \\
Années 5000 & Les terriens sont totalement enfermés dans leurs villes et ont peur des robots en général. Les planètes indépendantes (il y en a une cinquantaine) des \textit{Spacers} au contraire vont jusqu'à tester les robots humanoïdes (à peu près indiscernables d’un humain), mais finissent par abandonner l’idée. \\
Années 5200 & Un groupe de terriens se libèrent de l'hégémonie Spacer et entament une nouvelle phase de colonisation (l’idée des robots a été totalement abandonnée). \\
Années 10 000 & Les humains se sont répandus sur quasiment toute la Galaxie, et les luttes de pays d’autrefois recommencent à des échelles de distance beaucoup plus grandes. \\
Années 12 500 & La république de \textsc{Trantor} se crée : c’est l’an 1 de l’Ère Galactique. (on passera maintenant à ce calendrier pour la suite.) \\
Années 10 000 & C’est l’apogée de l’Empire Galactique de \textsc{Trantor} : il s’étend sur toute la Galaxie. \\
Années 12 050 & L’Empire commence de plus en plus à perdre le contrôle des sections de la périphérie de la Galaxie… Tels celui d’\textsc{Anacréon}. C’est ici que la Murder a lieu. \\
An 12 069 & An 1 de l’Ère de la Fondation. \\
\end{tabular}
Bon, je ne vous spoile pas trop la suite… Je gâcherais les \textit{Foundation}.

Un petit détail pour mettre de l’ambiance :
ci-dessous le logo représentant l’Empire. (présent sur tous les uniformes impériaux par exemple…)

\begin{center}
	\includegraphics[width=2cm]{logo_empire.png}
\end{center}

\section{Présentation générale de l’action}

\paragraph{À propos du système politique}
{
J’ai pris beaucoup de liberté sur ce que signifiait le terme «~gouverneur royal~» sur \textsc{Anacréon}.
En l’occurrence, j’ai choisi pour les besoins de l’intrigue le sens d’un gouverneur tout puissant, ou presque, mais qui est soumis à des «~votes~» réguliers pour savoir s’il reste en place.
Cependant, les gouverneurs restent «~royaux~» : pour pouvoir se présenter gouverneur à une élection, il faut avoir été choisi par une «~personne de confiance~» (un autre gouverneur ou un envoyé de \textsc{Trantor} par exemple).

J’espère que cette entorse au monde d’\textsc{Asimov} ne froissera personne.
}

\paragraph{Lieu de l’action}
{
On se place dans un univers «~fortement inspiré~» des \textsc{Foundations} d’Isaac \textsc{Asimov}, juste avant la colonisation de la planète \textsc{Terminus} par la première fondation. L’Empire existe encore, même s’il présente des signes de faiblesse assez marqué en province.

Toute l’action se passe dans un vaisseau spatial naviguant dans la province d’\textsc{Anacréon}, parmi les plus éloignées de la capitale de l’Empire, \textsc{Trantor}.

Ci-dessous un plan local du secteur et des secteurs voisins, avec les planètes importantes :
\begin{center}
	\includegraphics[width=10cm]{galactographie.png}
\end{center}
}

\paragraph{Synopsis}
{
Le gouverneur d’\textsc{Anacréon} et plusieurs de ses conseillers sont en train de se faire inspecter par un homme de l’Empire (venant de \textsc{Trantor}).
Le vaisseau se dirige vers \textsc{Anacréon} (il venait de sa lune habitée, où on a accueilli l'ambassadeur).
}

\paragraph{Avertissement}
{
Certains d'entre vous possèdent ou viendront à posséder des blasters ou d'autres armes. Comme vous incarnez des personnages normaux, il faut que vous ayez conscience qu'un SEUL tir est mortel si il vise certaines parties du corps (poitrine, tête,...). Réfléchissez donc bien avant de les utiliser.
}



\newpage

\mj{
\section{Présentation globale des personnages}

\paragraph{Paul \textsc{Ledar}}
{
Le gouverneur d’\textsc{Anacréon}.
Il détourne des montagnes d’argent qu’il utilise en grande partie en dessous de table pour se faire renommer gouverneur tous les quatre ans.
}

\paragraph{Leopold \textsc{Ledar}}
{
Le fils du gouverneur, qui aimerait bien avoir le pouvoir un jour.
Son pouvoir politique n’est pas forcément si important, mais la pression de son père sur la planète le fait croire très influent.
Il est officiellement parmi les conseillers du gouverneur (même si ce n’est pas de lui que sortira la moindre idée constructive) et c’est lui le préfet de \textsc{Meneleus}.
La seule chose dont il est courant pour les disparitions sur \textsc{Meneleus}, c’est que l’hypothèse d’un trafic d’armes et de soldats pour aller combattre sur \textsc{Ifni}, une planête du secteur voisin \textsc{Smyrno}, semble la plus probable vu les affrontements qu’il s’y produit actuellement (bien que le gouvernement de \textsc{Smyrno} tente d’étouffer au maximum l’affaire ; \textsc{Trantor} ne sait d’ailleurs absolument rien de tout cela\ldots).
L’autre piste possible est que ce soit la rébellion qui ait installé sa base ici et qui recrute activement.
}

\paragraph{Danrell \textsc{Sonen}}
{
Un des conseillers du gouverneur.
C’est en réalité un robot réactivé au tout début de l’Empire par un Spacer extrémiste.
Il n’a pas vraiment réfléchit pendant ses millénaires d’existence à une quelconque loi zéro.
Il cherche à tout moment à ne pas laisser paraître le moindre indice de sa non-humanité.
Son but est de faire respecter les lois (il n’en invente que très peu), quelque soit l’absurdité de celles-ci (du moment que l’on ne parle pas de tuer directement des êtres humains).
C’est ainsi que, bien qu’ayant suffisamment de détails pour savoir que Paul \textsc{Ledar} détourne de l’argent pour se faire réélire, il ignore tout simplement cela : ce n’est pas à lui de s’en occuper.
Mais lorsqu’il y réfléchit un peu plus, le côté légal de toute cette affaire semble pour le moins douteuse.
Il a entendu parler des réseaux de drogue à \textsc{Anacréon} et c’est l’affaire sur lequel il réfléchit activement actuellement.
}

\paragraph{Briran \textsc{Corllek}}
{
Un autre conseiller du gouverneur.
Celui-ci est beaucoup moins honnête et est le meneur principal d’un réseau d’esclavage (ces derniers sont enlevés à \textsc{Meneleus}) et de drogues (sur \textsc{Anacréon}).
Autant que possible, s’il pouvait étouffer les rumeurs venant de \textsc{Meneleus}, il le fera.
Il n’a jamais hésité à assassiner (jamais directement) des personnes qui en savaient trop.
Il profite beaucoup que Paul \textsc{Ledar} soit suffisamment naïf.
Il préfère par contre largement être celui qui tire les ficelles dans l’ombre à devenir gouverneur.
}

\paragraph{Vlad \textsc{Shyvock}}
{
C’est l'inspecteur.
S’il a choisi cette mission, c’est bien parce qu’il a à faire dans les environs.
Il connaît en effet l’existence du réseau d’esclavage de Briran \textsc{Corllek} en tant que\ldots client.
Il compte bien rentabiliser une petite planète inhabitée (car trop froide) contenant d’après les scanners de nombreux minerais précieux.
Il est donc favorable à ce petit «~commerce~» et empêchera quiconque fera mine de vouloir l’enlever.
Mais les affaires deviennent compliquées puisque les «~incidents~» sur \textsc{Meneleus} commencent à se faire connaître.
Il veut éviter à tout pris de devoir mettre ceci sur son rapport — à moins bien sûr que l’on trouve un faux coupable à arrêter.
On lui a de plus collé un fonctionnaire qui fouine un peu trop et il aimerait trouver moyen de s’en débarrasser.
}

\paragraph{Nuran \textsc{Hazon}}
{
Le «~compagnon~» de l’inspecteur.
Il a été envoyé par l’administration de \textsc{Trantor} sous la pression (dicrète néanmoins) de nombreux rivaux de Vlad \textsc{Shyvock}.
Il a comme mission officielle d’aider l’inspection, mais comme mission officieuse de trouver le plus possible d’éléments pouvant décrédibiliser \textsc{Shyvock}.
Il a actuellement besoin d’argent, et c’est la principale raison pour laquelle il a accepté la mission : il est devenu de plus en plus dépendant de drogues exotiques, et le fait de devoir se projeter en province ne l’enchante pas tellement. C’est un drogué, certes, mais pas régulier : il a des vagues d’humeurs, ses «~périodes~».
Si seulement il pouvait trouver une petite friandise locale pour se soulager un petit peu.
}

\paragraph{Plaïa \textsc{Idaski}}
{
C’est elle le chef de la sécurité.
Elle sympathise avec le groupe de résistance d’\textsc{Anacréon} contre le gouverneur toujours réélu dont les soupçons de détournement d’argent découverts commencent à devenir trop nombreux.
Cependant, elle ne veut en aucun cas perdre sa vie, ou la passer en prison pour le reste du temps.
Elle fera ce qu’elle peut pour qu’aucune confrontation se passe lors du transport (les caméras sont partout dans ce vaisseau, mais la plupart sont retransmises directement au central sur la planète mère).
Mais elle ne fera rien pour emprisonner les résistants du vaisseau, elle est d’ailleurs au courant pour le responsable des communications.
Elle doit briefer au tout début du jeu Dacan \textsc{Calrines}, qui semble étrange aujourd’hui.
}

\paragraph{Nuvan \textsc{Harsam}}
{
Le responsable des communications.
C’est un agent de la résistance locale assez engagé : il n’hésitera pas une seconde à mettre sa vie en danger si cela peut permettre de rétablir une démocratie ou quelque chose de ressemblant sur \textsc{Anacréon}.
Deux types de communications vont transiter : celles de la résistance et les autres.
La plupart des machines sont automatisées, il s’occupe juste de faire l’interface avec le reste du vaisseau (il a une puce électronique dans l’oreille qui lui indique du nouveau, ce qui lui permet de se déplacer librement dans le vaisseau).
Il apprendra au cours de la partie que parmi les personnes présentes sur le vaisseau, au moins l’une d’elle fait des affaires à \textsc{Anacréon} avec les réseaux de drogue ; découvrir son identité sera très utile à la rébellion, puis une fois découverte, l’éliminer pourra être un objectif (tout dépend si la rébellion pourra se servir de l’information).
Puis rebelote, mais pour des affaires à \textsc{Meneleus} (la rébellion ne sait pas exactement de quoi il s’agit).
}

\paragraph{Denik \textsc{Calkor}}
{
Garde du corps de Paul \textsc{Ledar}.
Ce qui l’intéresse, c’est l’argent, les médailles à mettre sur son bel uniforme et tout les avantages que l’on pourrait imaginer.
Extérieurement, il semble pourtant incapable d’accepter la moindre condition qui le détournerait de sa mission.
Il porte un bon blaster et il est capable de menacer ou de tuer quelqu’un pour obtenir des informations sans hésiter.
On lui laisse toujours la possibilité d’aller voir le gouverneur (qu’il discute ou non avec un autre personnage important) et en connaît donc parfois long.
Il est censé obéir au gouverneur sans discuter. (Sauf si celui-ci lui demande de partir.)
Les personnes présentes sur le vaisseau ont l’air louches, et il aimerait en découvrir plus, puis — suivant le cas — en soutirer de l’argent contre la promesse de ne pas divulguer d’informations, ou les éliminer (d’une manière ou d’une autre : on peut tout simplement les dénoncer parfois\ldots).
Il est un client assez régulier du réseau de drogue locale, heureusement que personne ne penserait à faire des tests très poussés sur sa personne (ce serait plutôt à lui d’exiger de quelqu’un un tel contrôle.).
Dacan \textsc{Calrines} semble avoir un petit stocke chez lui, mais il n’est en pas sûr.
Il va chercher à faire affaire avec.
}

\paragraph{Avalga \textsc{Orzon}}
{
Une très riche propriétaire d’\textsc{Anacréon}.
C’est aussi une amie politique du gouverneur.
Il y a déjà eu de nombreux échanges d’argent dans les deux sens entre eux (rarement légal).
Elle a déjà eu de nombreux problèmes avec la résistance, qui se doutait que les jugements étaient truqués.
Très intéressée par l’argent ou tout autre avantage, elle marchande tout ce qu’elle peut.
Elle va chercher à créer de nouveaux accords commerciaux avec les différents personnages importants du vaisseau, honnêtes ou non, du moment qu’ils rapportent.
Elle a entendue parler par exemple d’un traffic d’armes sur \textsc{Meneleus}, et ça l’intéresse grandement.
}

\paragraph{Yovan \textsc{Bunski}}
{
Le médecin du vaisseau.
Il n’a jamais porté dans son cœur le gouverneur, mais il reçoit du fait de sa position tant de menaces sur sa famille qu’il ne fera pas grand chose qui risquerait de l’inculper.
Néanmoins, il sympathise avec la rébellion et fera tout ce qu’il peut pour les aider tant que sa famille ne risquera rien (c’est à dire que l’on ne puisse pas savoir qu’il est impliqué là dedans).
Il a découvert dans des toilettes collectives une poudre qu’il a analysé comme étant de la drogue similaire à celle du traffic sur \textsc{Anacréon}.
Il ne dira rien à personne sans savoir qui ici est un drogué, mais ce genre d’informations peuvent être très utiles.
}

\paragraph{Vanri \textsc{Thrawder}}
{
Espionne du secteur d’à côté, \textsc{Smyrno} (horiginaire de la planête principale, \textsc{Smyrno}).
Elle cherche à enquêter sur le gouverneur et les différentes magouilles qu’il s’y passe par ici.
Elle connaît la rébellion de réputation et aimerait se faire connaître en tant qu’allié d’eux (oui, elle joue triple jeu).
Son but est de faire couler le secteur d’\textsc{Anacréon} et fera tout pour que la situation politique ici soit le plus instable possible.
C’est elle qui a remplacé les pilules contre la toux de Dacan \textsc{Calrines}, la seule personne un peu honnête ici, par des pilules mnémosiques et a convaincu plusieurs personnes (en laissant des indices, pas en leur disant) qu’il ferait un traffic d’armes dans la région de \textsc{Meneleus}.
Elle se fait passer pour la responsable technique, ce qui lui a permis d’infiltrer quelques mouchard qui pourront être récupérés pendant la partie.
Personne ne lui fera la moindre remarque sur son travail : c’est quasiment la seule ici qui s’y connaît un minimum en technologie.
}

\paragraph{Dacan \textsc{Calrines}}
{
Il ne se rappelle même pas de son nom, à par en lisant la plaque inscrite sur son uniforme.
En fouillant dans ses affaires, il a découvert une caisse de blaster (déposée par Vanri \textsc{Thrawder}), mais il n’en sait rien, ainsi qu’un tas non négligeable de crédits.
Parmi les documents qu’il lit, il semblerait qu’il soit à la tête d’un traffic d’armes sur \textsc{Meneleus}.
Bien entendu, rien ne lui revient en mémoire.
Il remarque un mot de lui dans sa chambre : il a rendez-vous avec Plaïa \textsc{Idaski} pour un briefing de la mission actuelle.
Avant que Vanri \textsc{Thrawder} mette le désordre dans sa tête, il était juste un agent de sécurité aux ordre de Plaïa \textsc{Idaski}.
Pas des plus finaux, mais il faisait toujours bien son travail.
Il était un des nombreux «~terminaux~» du traffic de drogue.
Il n’a jamais fait de nombreuses affaires, mais il a acheté des quantitées relativement grandes de drogue (pour une dizaine de personnes), qu’il comptait revendre à \textsc{Anacréon}.
Cette drogue est cachée dans un oreiller, et il l’a repéré lors de ses fouilles.
Par contre, il ne voit pas du tout ce que ça peut bien être\ldots
}

\newpage
}

\player{1}{
\section{Ton personnage : Paul \textsc{Ledar}}

\paragraph{Description physique}
{
	Pas spécialement grand, en dessous de la moyenne anacreonnienne.
	Il est cependant très bien habillé (c’est le gouverneur).
}

\paragraph{Ton personnage}
{
Toi au moins, tu as bien compris ce que c’était que la politique !
Cela fait un certain temps que tu t’y es mis. Jamais tu n’a essayé autre chose.
La politique, c’est ta vie. Jamais tu ne laissera quelqu’un te marcher sur les pieds.

Cela va bientôt faire dix ans que tu occupe le poste de gouverneur d’\textsc{Anacréon}.
Gouverneur, c’est une bonne position, surtout d’un système aussi éloigné du centre de la Galaxie, \textsc{Trantor}.
Habituellement, un gouverneur est réélu tous les quatres ans. Mais bon… Les dessous de tables ont toujours existé.
Et tu n’as aucune intention de lâcher ce poste.

Tu as déjà réflechi à la question de supprimer cette loi qui fait apparaître un semblant de démocratie à ta dictature, mais tu penses que jamais le peuple anacreonnien n’acceptera ceci.
Cela t’embêtes beaucoup d’ailleurs : cela te coûte de plus en plus cher.
Certes, tu détourne plus d’argent que tu en utilise pour te faire élire, mais ça t’embête tout de même.

Heureusement qu’il y a ton garde du corps Denik \textsc{Calkor}.
Ce serait un suicide de s’en débarrasser plus d’un quart d’heure, surtout par les temps qui court.
Un grand gars bien bourru qui t’obéit au doigt et à l’œil, c’est toujours pratique.

Tu connais bien Avalga \textsc{Orzon}, c’est une riche propriétaire avec qui tu as souvent fait affaire.
Il serrait de bon goût pour toi de conclure une nouvelle affaire ici : si seulement la moitiée de ses relations étaient à ton service pour truquer les élections, ce serait tellement plus simple.
Il faudrait lui donner quelques terres sur une planète du système.

En parlant de petite planète, ça te fait penser à ton fils Leopold \textsc{Ledar}.
Quel boulet, lui, il n’a vraiment rien compris à la vie.
Tu ne sais pas ce qu’il s’imagine, mais il serait bien de lui faire fermer sa grande gueule.
C’est pas pour dire, mais là, il semble devenir vraiment gênant.
Il est en effet préfet de la planète la plus petite du système, \textsc{Meneleus}.
Pour l’instant, la situation politique de la planête est assez instable paraît-il…
Et si tu donnait la préfecture de la planète à Avalga \textsc{Orzon} ?
Mmmh… Il te faudrait une excuse : tirer une balle dans sa propre famille, cela risque de ne pas plaîre à ton électorat.

Parlons-en de ton électorat d’ailleurs !
Depuis cette fichu rébellion, il ne cesse de descendre !
Mais tu ne sais pas vraiment où peuvent se trouver tous ces rebelles.
Il faudrait que tu demandes un peu d’aide à tes fidèles conseillers Danrell \textsc{Sonen} et Briran \textsc{Corllek}, qui sont assez bon conseillers, il faut le dire.
(Surtout Briran \textsc{Corllek}, qui a déjà montré une excellente intuition des réaction violentes que le peuple d’\textsc{Anacréon} peut parfois avoir.)
Ton fils aussi est conseiller, il fallait bien lui trouver un poste pas trop loin pour le surveiller, mais bon, voilà.

Le vaisseau est assez grand et tu ne connais pas tout le monde.
Serais-ce possible que des rebelles se cachent ici même ?
Non, ils auraient dus passer les sécurités, et Plaïa \textsc{Idaski} est assez ferme sur ce genre de points.
Mais bon, heureusement, tu as ton garde du corps.

Ce vaisseau est actuellement en route pour \textsc{Anacréon}, la planète principale du système.
Il vient d’acceuillir l’ambassadeur de l’Empire.
Ça, c’est vraiment ton plus gros problème actuellement.
Tu sais très bien que tout ne se passe pas très bien dans le système, et tu n’as aucune envie de devoir faire un rapport à un représentant de \textsc{Trantor}.
Mais bon, il n’y a aucune solution : on ne va pas chercher des ennuis avec l’Empire.

Par contre, cela ne veut pas dire du tout que tu vas leur annoncer la vérité sur un planteau.
Non, il va falloir improviser, mais il est hors de question que \textsc{Trantor} se doute de quoi que ce soit sur ce qui ce passe exactement ici.
}
}

\player{2}{
\section{Ton personnage : Leopold \textsc{Ledar}}

Tu es le fils du gouverneur Paul \textsc{Ledar}.
Ta vie, ça été très simple à chaque fois : en tant que fils d’un homme politique influent, puis du gouverneur, tout de tombait sous la main.

D’abord les études. On a essayé de te faire gober des montagnes de connaissance en économie.
Mais bon, des élèves sérieux, ça ne court pas les rues ici…

Et puis toi, d’abord, tu commence à en avoir marre d’être traité comme le fils à Papa !
Si seulement tu pouvais être vraiment gouverneur maintenant.
Là tu pourrais faire taire tous ces gens !
Ils t’énervent ! D’accord, tu n’a pas beaucoup souffert dans la vie pour trouver ce poste de conseiller du gouverneur, mais ce n’est pas une raison !

Ah ! Si seulement tu pouvais être gouverneur…

À vrai dire, le poste de conseiller te plaît aussi, c’est juste que tu sens que personne ne t’écoute.
Tu est conseiller, mais Paul \textsc{Ledar} n’écoute que les autres conseillers Danrell \textsc{Sonen} et Briran \textsc{Corllek} !
C’est comme s’il ne te faisais pas confiance, c’est énervant !

Et puis cette panète dont ton père t’a donné la préfecture, comment s’appelle-t-elle déjà ?
Ah oui, \textsc{Meneleus}.
Une vraie blague.
C’est une planète pourrie, isolée de tout dans le système d’\textsc{Anacréon}.
Toi, c’est \textsc{Anacréon} elle-même que tu voulais, pas cette planète de bouseux complètement gelée !

Heureusement, cette position te donne quelques pouvoirs, et tu les utilises de temps en temps — il faut bien s’amuser !

Actuellement, tes contacts sur la planète d’indiquent que des disparitions régulières semblent arriver…
Bah ! Ce n’est pas ces quelques bouseux qui vont t’embêter !
Qu’ils disparaissent s’ils le veulent.
On ne cesse cependant de le te répéter…
Certains parlent d’un traffic d’armes et de soldats pour alimenter la guerre sur \textsc{Ifni}.
Il s’y produit en effet ces temps-ci des affrontements assez beliqueux paraît-il.
Mais tu n’en sais pas plus : le secteur de \textsc{Smyrno} où se trouve \textsc{Ifni} semble vouloir étouffer l’affaire.
Il paraît que même \textsc{Trantor} n’est pas au courant !

Bah ! \textsc{Ifni} est loin (le secteur d’à côté quoi !), de quoi peuvent-ils se mêler ?
Qu’ils crèvent là-bas, ça te va très bien.

Par contre, une autre piste serait que la rébellion s’est installée sur la planète et qu’elle recrute activement.
Qu’est-ce qu’ils sont bêtes ! Il seraient faient comme des rats : Tu n’a qu’à demander la destruction de la planète et se sera réglé.

En plus, cela te permettra de vaincre la rébellion.
Si cela ne te promeut pas gouverneur, tu ne vois pas ce qu’il faut !

Dommage qu’il faille l’accord de ton père le gouverneur pour détruire une planète du système, ça aurait été réglé depuis lontemps.
}

\player{3}{
\section{Ton personnage : Danrell \textsc{Sonen}}

Ton personnage est en réalité un robot :
Avant que l’empire n’existe encore, il y avait deux «~clans~» — les \textit{Spacers} et les \textit{Settlers}.
Pour faire vite : les Spacers vivaient sur une dizaine de planêtes, pouvaient vivre jusqu’à plusieurs centaines d’années, et vivaient au milieu d’une foule de robot dans une société où le taux de natalité était contrôlé afin d’éviter toute surpopulation. Ils avaient eu pendant très longtemps le contrôle politique de la Terre.
Les Settlers quant à eux sont d’anciens Terriens qui ont réussi à se libérer de l’hégémonie des Spacers. Ils refusaient tous robots, et ont permis de rapidement construire ce que l’on appelle maintenant l’«~Empire~».

Tu as été réactivé bien après la fin de l’hégémonie Spacer par un Spacer extrémiste voulant donner une chance aux robots.
Comme tout robot, tu suis les trois lois de la robotique :
\begin{itemize}
	\item Jamais tu ne fera du mal à un humain, que se soit en faisant une action ou en ne faisant pas une action.
	\item Tant que les ordres des autres humains ne contredisent pas la règle précédente, tu obéit aveuglement aux humains.
	\item Tant que tu ne désobéit pas aux deux lois précédentes, tu tente autant que tu peux de te protéger toi-même.
\end{itemize}

D’apparence extérieure, tu ressembles totalement à un humain, aux différences près que tu as tendance à paraître sans émotions quasiment tout le temps.
Par contre, comme tout robot, tu as ue force extra-ordinaire. (Mais tu ne la montre jamais, sauf si cela contredit les deux premières lois.)
Contraierment à Daneel \textsc{Olivaw}, tu n’a jamais vraiment réfléchi à la possibilité d’une «~loi zéro~».

En tant que robot, tu as bien vu que si tu restait à \textsc{Trantor}, le centre de l’Empire, jamais tu ne pourrais éviter d’être repéré comme étant un robot (ne serais-ce que parce que tu ne veilli pas).
Tu as donc décidé (suivant la troisième loi) de partir dans les périphéries de la Galaxie.

Cela va bientôt faire près de 13 000 ans que tu as été activé.
Depuis ce temps là, tu est passé de systèmes en systèmes dans les périphéries de la Galaxie.
Tu te fait passer pour un humain, tu tentes de te faire le moins remarquer possible, tout en prenant part à la vie humaine.
Assez souvent, tu restes environ 60 ans sur une planête avant de changer de système.

Cette fois-ci, il semble que tu ait été particulièrement chanceux :
En arrivant sur \textsc{Anacréon}, tu t’es très vite fait repéré en sauvant un habitant.
Tu t’es assez vite lié d’amitié avec lui. (Depuis le temps, tu as l’habitude de parler sans jamais te trahir.)
Ta capacité à calculer des probabilités d’actions ou de décision c’est assez vite trouvée utile ici.
Tu as réussi à te faire nommer conseiller du gouverneur.

Comme tout robot, ton idée de la «~Justice~» est assez simpliste : suivre les lois (de la robotiques, comme les lois humaines).
Inventer une loi n’est pas ton fort, sauf si cela renforce une autre.
Souvent, tu te contente d’indiquer au gouverneur Paul \textsc{Ledar} ce que tu considère comme étant l’action politique la plus probable.

Tu as déjà eu suffisament d’indices pour voir que Paul \textsc{Ledar} détournait des quantités d’argent dans le seul but de se faire réélire.
Tu trouve cela assez louche, mais comme certaines lois sont ambigües, tu ne voit vraiment le problème.
En fin de compte, tu ignores simplement cette affaire, même si tu pourrais la lâcher sans le faire exprès (mais jamais de manière volontaire, sauf si on t’a demandé ce que tu savais sur lui).

Tu sais très bien ce qui pourrais se passer si on te donnait un ordre : tu le suivrais sans même réfléchir.
Pour éviter cela, tu as appris à ne jamais te méler d’affaires qui ne te conscerne pas, à rester proche de Paul \textsc{Ledar}, où jamais personne n’oserait te donner un ordre directe. (Tu est tout de même le conseiller du gouverneur !)
Depuis le temps que tu est activé, tu as su trouver des «~limites~» aux lois robotiques : si quelqu’un te demande ce que tu sais sur tel sujet, il ne te demande pas directement de dire tout ce que tu sais, mais juste de dire ce que tu pense \emph{pourrais} être vrai compte-tenu de tes informations par exemple.
Tu as su jouer sur ce tranchant de la langue pendant un certain temps.

Actuellement, ton rôle de conseiller est de mettre au clair ce qui se passe exactement à \textsc{Anacréon} conscernant les traffics de drogue.
Pour l’instant, la seule chose que tu sais dessus est la précence de quelques réseaux.
En souppesant les probabilités, tu penses que la majeur partie se trouve dans la ville de \textsc{Concatalle}, mais tu n’en sais pas encore assez pour pouvoir effectuer la moindre action de la part du gouvernement.

\paragraph{Capacités spéciales}
{
Ton personnage possède une force surhumaine.
Il a aussi une très grande connaissance d’énormément d’affaires louches dans la Galaxie.

Bien entendu, il ne va pas le montrer à tout le monde, de peur de montrer aux autres sa véritable identité ; mais s’il se produit un cas où il se retrouve avoir besoin de le faire (par exemple pour obéir à une des lois), merci d’appeler un maître du jeu.
}
}

\player{4}{
\section{Ton personnage : Briran \textsc{Corllek}}

Dans ta vie, tu as très rapidement compris ceci :
les gens sont faibles et facilements manipulables à ta volontée.
Beaucoup sont cupides et une simple somme d’argent peut suffire à les faire changer d’avis, de parti politique ou même à les faire oublier quelque chose.

Depuis que tu as compris ceci, rien ne t’arrête.
Tu désir depuis ce temps là d’avoir le plus possible de pouvoir.
Tu sais manipuler les lois de façon à leur faire dire ce qe tu veux, tu n’as jamais hésité à soudoyer, menacer, torturer ou faire tuer pour arriver à tes fins.

Tu es relativement rapidement devenu un des meneurs principaux des réseaux d’esclavage et de drogue du secteur.
Le réseau d’esclavage est un réseau enlevant des travailleurs sur une petite planète froide du secteur, \textsc{Meneleus}, pour les exporter dans d’autres secteurs de la Galaxie.
Quant au réseau de drogue, tout se passe sur \textsc{Anacréon}.
La drogue locale est une poudre brunâtre qui se vend assez cher.

Bien entendu, rien de tout cela n’est légal, mais bon, pourquoi suivre les lois si cela reviendrait à perdre une partie de ses pouvoirs ?
Tu n’es pas un de ces faibles !

Tu as de plus réussi à avoir une des meilleurs positions auxquelles tu pouvais espérer :
Tu es un des conseillers du gouverneur Paul \textsc{Ledar}.
Paul \textsc{Ledar}, l’exemple typique d’un humain \textit{faible}.
Il suit quasiment tous tes «~conseil~» à la lettre — bien entendu, tu fais tout ce qu’il est possible de faire pour que ces conseils parraissent à son avantage.

Tu as vite découvert que suivre les lois n’était pas non plus dans ses objectifs, et tu l’a fait lancer une grande campagne de dessous de tables pour rester au pouvoir.
Comme s’il était le premier !
Ce qui est rigolo avec lui, c’est qu’il a des… remords.
Quel faible !

En pratique, tu es donc la personne qui tire les ficelles derrière l’ombre, qui contrôle tout le gouvernement d’\textsc{Anacréon}.
Ça te va très bien.

Bien entendu, il est déjà arrivé de nombreux cas où il n’a pas suivi tes conseils.
Il l’a vite regreté : à chaque fois, tu faisais assassiner un de ses proches et faisait passer cet assassinat comme une réponse du peuple à son innaction.
Et il te croit !

Il y a actuellement quelques rumeurs venant de \textsc{Meneleus} à propos de ton traffic d’esclaves.
Si seulement tu pouvais l’étouffer.

Il y a de plus cet envoyé de \textsc{Trantor} qui peut tout gâcher.
Quels fouineurs à \textsc{Trantor}, comme s’ils avaient encore du pouvoir ici !
Comment vas-tu faire pour t’en débarrasser ?

\paragraph{Compétences spéciales}
{
Tu connais très bien les réseaux d’assassinat d’\textsc{Anacréon}, et vous avez des «~codes~» qui vous permettront de communiquer sans que personne ne comprennent.
Si tu souhaites assassiner un proche de quelqu’un dans ce vaisseau, tu n’as qu’à prévenir un des maîtres du jeu.
(Attention, l’assassinat pourra prendre quelques temps avant d’avoir effectivement lieu !)
}
}

\player{5}{
\section{Ton personnage : Vlad \textsc{Shyvock}}

Tu viens de la planète \textsc{Trantor}, le centre administratif de l’Empire Galatique.
Pourquoi te retrouver ainsi perdu au milieu d’un secteur de la périphérie de la Galaxie ?
C’est une longue histoire en fait…

Officiellement, tu es un fonctionnaire de \textsc{Trantor}.
En réalité, ta principale «~cible~» actuelle n’est pas ces paperrasses de fonctionnaires, non.
Tu as récement «~acheté~» le monopole d’une planète relativement récement découverte (une dizaine d’années).
Il va sans dire qu’«~acheter~» est synonyme ici de dessous de tables, de corruption, etc.

Cette planète totalement inhabitée (car beaucoup trop froide) contiendrait d’après les scanners de nombreux stocks de minerais précieux.
Tu cherches donc de la main d’œuvre pas cher capable de s’adapter à une telle planète.

Et une telle main d’œuvre, tu sais où en trouver ici.
En effet, tu connais de sources à peu près fiables qu’un certain Briran \textsc{Corllek} organise un traffic d’esclaves venant d’une planète froide nommée \textsc{Mélènus}, ou quelque chose comme ça.
Cette personne a l’air d’avoir pas mal d’antécédants : tu es sûr de pouvoir faire affaire avec.

C’est la raison principale pour laquelle tu as accepté cette mission d’inspection.
Ta mission officielle est de faire un rapport exact sur ce qui se passe dans la périphérie, en particulier tout ce qui pourrait coller avec toute contradiction avec les lois officielles de \textsc{Trantor}.

Bien entendu, il serait de bon ton pour tes affaires de ne pas mettre par écrit cette histoire d’esclavage… À moins bien sûr de trouver des «~preuves~» contre un faux coupable qui serait arrêté par le gouverneur local, Paul \textsc{Ledar}.
Il faudrait aussi conclure ce pacte avec ce \textsc{Corllek}.
Tu as vraiment besoin de tels esclaves.

Ah oui, il reste un petit détail.
On t’a collé un fonctionnaire pour t’assister dans ta mission officielle et celui-ci a une fâcheuse tendance à fouiner un peu partout — un peu trop aussi.
Si seulement tu pouvais t’en débarrasser… Ça te simplifierait de beaucoup la tâche : s’il en venait à découvrir cette histoire de traffic d’esclave, tu serrais forcé à le mettre sur ton rapport tel quel.
}

\player{6}{
\section{Ton personnage : Nuran \textsc{Hazon}}

Tu vient de la planète \textsc{Trantor}.
Tu y occupe un rôle de fonctionnaire de l’administration de la Capitale Impériale.
Si tu es ici, c’est officiellement pour assister Vlad \textsc{Shyvock} dans une mission ici.
Ce dernier doit en effet faire une mission de reconnaissance dans le secteur d’\textsc{Anacréon}, il paraîtrait que des choses louches s’y passent.
C’est une région de la périphérie de la Galaxie qui semble perdre un peu les derniers fondements de démocratie qui s’y appliquait.
Bon, ça, tu t’en fiche. Que ces hommes qui ne respectent pas les lois de l’Empire se débrouillent comme ils veulent, du moment qu’ils ne se rapprochent pas du centre.
(Mais ça, tu n’y croit juste pas : qui oserait affronter l’Empire de nos jours ?)

La vraie raison de ta précense est toute autre.
Vlad \textsc{Shyvock} dérange pas mal les affaires de l’administration trantorienne, et plusieurs de ses opposants montrent une vive insistance pour qu’il soit… mis hors fonctions.
Cette pression est certes discrète, elle reste perceptible, et tu as réussi à négocier le prix fort à l’un de ses rivaux pour l’accompagner jusqu’ici.

C’est assez dangereux.
Tu pourrais y perdre ton poste à \textsc{Trantor} si tu te fait prendre.
Mais cela ne devrais pas arriver, tu sais faire les choses subtiles.
Il te faut donc trouver tout ce qui pourrait nuir ou décrédibiliser \textsc{Shyvock} à \textsc{Trantor}.
Tout peut servir.
Tu as pour cette mission acheté un petit appareil photographique.
Chaque photo «~intéressante~» peut te rapporter très gros.

La raison pour laquelle tu as besoin d’argent est toute autre.
Depuis un certains temps, tu te drogues régulièrement.
Certes, pas toujours, tu as tes «~périodes~».
Ce que tu préfère, ce sont les drogues «~exotiques~».

Certes, d’un point de vue politique, la périphérie de la Galaxie est tout sauf intéressante, d’un point de vue «~exotique~», elle regorge de surprise.
L’importation coûte cher, surtout pour ce type de marchandises.
Si seulement tu pouvais profiter du séjour pour obtenir quelques petites «~friandises~» locales, ça te soulagerait peut-être.

\paragraph{Matériel}
{
Ton personnage possède un appareil photographique miniature.

En cas d’utilisation, merci de prévenir un maître du jeu.
}
}

\player{7}{
\section{Ton personnage : Plaïa \textsc{Idaski}}

\paragraph{Description Physique}
{
Assez grande et costaude, genre garçon manqué.
}

Les détournements d’argent du gouverneur Paul \textsc{Ledar} commencent à bien faire, décidement, la résistance à raison : il faut en finir avec son gouvernement de dessous de table et de corruption !
Par contre, il faut le faire de manière subtile : l’idée même que tu pourrais y perdre la vie ou la passer en prison ne t’enchante guère…
Non, décidément, ce n’est pas toi qui tentera une quelconque action « héroique » ou suicidaire contre lui.
Par contre, rien ne t’empêche d’aider des partisans ou des membres de la résistance si l’occasion se présentait.
À condition bien sûr que tu ne mette pas tes mains là dedans : hors de question de se faire prendre sur le vif !

Tu penses d’ailleurs que tu va avoir un quelques occasions lors du voyage actuel.
En effet, c’est toi le chef de la sécurité du vaisseau menant le gouverneur vers la planète \textsc{Anacréon}.
Il revient de la lune où il est allé souhaiter la bienvenue à un ambassadeur de \textsc{Trantor}, qui l’accompagne d’ailleurs actuellement.

Bon, tout ça ne te regarde pas.
Ce qui est important pour toi, c’est tout d’abord d’assurer la sécurité dans ce vaisseau.
Après, si l’occasion se présentait de pouvoir tuer le gouverneur ou renverser le gouvernement — et que tu ne risque absolument rien — tu feras probablement quelque chose.
Mais tu n’es pas encore là.
Pour l’instant, la sécurité du vaisseau d’abord.

Et puis de toute façon, que ferais-tu ?
Il y a des caméras partout dans le vaisseau, et la plupart retransmettent directement vers la planète mère.
Inutile de se faire des histoires : si tu tire dans la direction du gouverneur, tout le monde saura que c’est toi une fois le vaisseau atteri.
Et tu ne le veux pour rien au monde.

Par contre, il pourrait être intéressant de voir ce que fera le soit-disant responsable des communications, Nuvan \textsc{Harsam}.
En effet, tu es quasiment certaine que cet homme fait partie de la résistance.
Après tout, s’il te «~vole~» ton blaster, comment pourrait-ce être de ta faute ?
Il reste à discuter discrêtement avec lui pour voir ce qui se passe.
Remarque, en tant que chef de la sécurité, ça doit pouvoir s’arranger.

Certes cette position donne quelques avantages, mais elle donne aussi quelques responsabilités.
Tu connais les actions à faire en situation d’urgence par cœur, et ce sera à toi de les rappeler à tout le monde en cas de problème (qu’il soit dû à un évènement extérieur ou à un tué/blessé dans le vaisseau).

Bien entendu, tu n’es pas toute seule à protéger ce vaisseau : Dacan \textsc{Calrines} est à tes ordres.
Il n’est pas spécialement du genre brutal, et c’est toujours agréable de travailler avec quelqu‘un comme ça.
Par contre, comme d’habitude, il va avoir oublié toutes les consignes que tu lui as donné la veille.

Bon, tu le vois arriver ; on peut dire qu’il n’est pas frais ce matin…
C’est parti pour une séance de rapide briefing histoire de lui rappeler comment tu as organisé la mission.

\paragraph{Matériel}
{
Ton personnage possède un blaster.
Libre à toi de l’utiliser comme tu le souhaite, à condition de prévenir un maître du jeu à chaque utilisation.
}
}

\player{8}{
\section{Ton personnage : Nuvan \textsc{Harsam}}

\paragraph{Description physique}
{
Assez jeune, pas très grand, plutôt maigre.
}

\paragraph{Ton personnage}
{
Depuis que est devenu majeur, tu comprend mieux combien le gouverneur Paul \textsc{Ledar} est un vrai connard.
Il n’a jamais eu aucune conviction politique, aucune action qui aurait permis d’aider la société du système.
Non, rien ! Par contre la corruption, ça il connaît bien.

Maintenant tu vois bien le système tel qu’il est : une dictature.
Certes, il y a encore des symboles (on offre une possible réélection du gouverneur), mais rien qui pourrait vraiment empêcher tout cela.

Il est maintenant temps d’en finir avec ce système qui a montré qu’il avait touché ses limites.
Il est temps de rétablir un vrai système démocratique comme autrefois.
Mais le seul moyen d’arriver à ces fins est la force, malheureusement.

Tu fait maintenant partie de la «~Rébellion~», le mouvement de résistance contre le gouverneur.
C’est un groupe secret qui a réussit à prendre une relative importance pour un tel mouvement.
Tu as déjà participé à des «~coups~» et a failli y perdre la peau.
Mais c’est pour la bonne cause : jamais tu n’hésitera une seule seconde à sacrifier ta vie si cela peut permettre de rétablir à coup sûr une démocratie ou quelque chose approchant.
La plupart de ces actions se faisaient dans la ville où les partisants pour la Rébellion sont les plus nombreux : la ville de \textsc{Concatalle}.

Mais là… là, tu es sur un coup vraiment très important.
Tu as su prouvé ta loyauté, toute la Rébellion te fait confiance pour cette mission.
Tu ne les décevra pas.

Tu as réussi à te faire embaucher dans un vaisseau du gouvernement, en tant que responsable des communications avec l’extérieur.
En pratique, c’est un petit rôle : la plupart des informations transites directement des ordinateurs de bord vers un central basé sur la planète d’\textsc{Anacréon}.
Ton rôle dans le vaisseau consiste à communiquer aux membres de l’équipage les informations importantes, ou de faire la requête inverse.

Mais ce qui est important est que des agents bien placés de la Rébellion on fait que le vaisseau sur lequel tu travaille depuis un certain temps sera celui qui transportera le gouverneur sur la planète depuis la lune de celle-ci.
Un ambassadeur de \textsc{Trantor} vient en effet d’atterrir sur cette dernière, et il est de bon ton pour un ambassadeur de faire bonne figure en l’accompagnant personnelement.
Ta mission est ainsi de faire tout ce que tu pourra pour assassiner Paul \textsc{Ledar}.

Tous tes faits et gestes sont probablement surveillés sur la planète (il y a des caméras quasiment partout sur le vaisseau, qui restransmettent probablement au central), inutile de ramener une arme (ou du poison ou n’importe quoi) ou d’en voler sur place, on ne te laissera pas faire et tu sera abattu avant même d’avoir pu faire quoi que ce soit.
Ta mort avant que tu ait fait ta mission sera une très grande perte pour la Rébellion : elle a travaillé très activement, et plusieurs hommes sont déjà morts pour que la situation actuelle soit ainsi.
Il s’agit donc de faire les choses bien.

Lorsque tu as enfin vu dans le vaisseau le gouverneur, il était cependant accompagné (outre les conseillers, et tout ce qui va avec) par un garde du corps qui avait l’air assez… inflexible.
Ça risque d’être coton.

Heureusement, des membres de la Rébellion ont réussi à trouver la fréquence de communication du vaisseau : peut-être la Rébellion en saura plus plus tard.

En pratique, actuellement, la mission est de trouver qui ici pourraient t’aider ; puis de voir si on peut passer à l’action.
Il ne faudrait pas perdre cette occasion : si lors de l’arrivée dans l’atmoshpère basse d’\textsc{Anacréon}, tu n’est toujours pas passé à l’action, il faudra prendre un peu plus de risques (sauf si ce garde du corps est trop efficace bien entendu, il ne s’agit pas de suicider).
}

\paragraph{Matériel}
{
Nuvan \textsc{Harsam} possède une oreillette.
Attention, il ne possède pas d’arme.

Lorsque le vaisseau reçoit une transmission de l’extérieur, les maîtres du jeu devront te prévenir.
}
}

\player{9}{
\section{Ton personnage : Denik \textsc{Calkor}}

\paragraph{Description physique}
{
Habits de garde du corps.
Extérieurement, Denik \textsc{Calkor} a l’air totalement impassible et incorruptible. (d’apparence…)

Il possède sur lui de nombreuses médailles pour le mérite.
}

Ton physique imposant et ta volonté à porter un bel uniforme t’as poussé à exercer un métier risqué.
Après avoir fait l’armée, où tu t’es efforcé à collectionner les médailles et autres signes extérieurs de ta supériorité.
Cette volonter d’avoir un bel uniforme et de belles médailles, tu l’as toujours.

Lorsque plus tard l’occasion s’est présenté, tu as choisi garde du corps.
Puis le gouverneur Paul \textsc{Ledar} lui-même t’a embauché !
En voilà du prestige ! c’est exactement ce qu’il te fallait.

Cette place, tu ne veux pas la risquer.
Ainsi, tu as toujours ton blaster sur toi et jamais tu n’hésitera à menacer ou tuer quelqu’un pour obtenir des informations utiles à ta mission de protection du gouverneur.
De plus, tu n’hésiteras pas à obéir aux éventuels ordres du gouverneur.

Tu as eu du mal à l’imposer au gouverneur, mais maintenant il ne faut plus une seule objection à ta présence :
jamais tu n’a de problème pour rester à côté de celui-ci ; même s’il est en pleine négociation secrète avec un personnage très important.

Tu en a déjà entendu pas mal de la part du gouverneur.
Notamment ses dessous de tables assez… fréquents pour ses multiples réélections.
Bon, c’est une chose. A priori, tu n’iras pas trop le répéter, tu tiens à ta place.

Remarque ta place… Depuis que tu a été accepté garde du corps personnel du gouverneur, tu t’es mis à la drogue d’\textsc{Anacréon}.
Elle se présente sous la forme d’une poudre brune, et elle est assez bonne pour toi.
Tu es un client relativement régulier du réseau local.
Heureusement que l’idée ne viendrait à personne de tenter des tests contre toi, se serait même plutôt à toi d’exiger un tel contrôle à quelqu’un si nécessaire.
(Le médecin de bord est censé être capable d’effectuer cette tâche, à tout hasard.)

D’ailleurs Dacan \textsc{Calrines} semble avoir un petit stock chez lui, mais tu n’es pas sûr du tout.
Il va falloir vérifier ça : tu es sûr que vous pouvez faire affaire sur ce coup là.

Quant à ta mission actuelle, tu sens qu’il va falloir faire attention :
les personnes présentes sur le vaisseau ont l’air vraiment louches.
Tu aimerait bien en découvrir plus.
Ensuite — selon le cas — ou alors soutirer de l’argent contre une promesse de ne rien divulger, ou les éliminer.

\paragraph{Matériel}
{
Tu porte sur toi un blaster.
Merci de prévenir les maîtres du jeu à chaque utilisation.
}
}

\player{10}{
\section{Ton personnage : Avalga \textsc{Orzon}}

\paragraph{Description physique}
{
Habillée d’une façon à montrer toutes ses richesses.
}

Depuis qu’il t’a aidé à accumuler toute cette fortune, toi et Paul \textsc{Ledar} s’aprécient particulièrement bien.
Pas trop de sous-entendus, vous êtes des gens fortunés qui ne pensent qu’à leur argent à longueur de journée… Mais tu aimes cette vie.
Comment pourrais-tu vivre sans ces richesses ?

Pour rien au monde du ne voudrais la perdre.
Pour cela, tu as toujours eu le sens des affaires ; et conclure un marché avec une autre personne est toujours bon pour les affaires justement.

C’est vraiment Paul \textsc{Ledar} l’homme avec qui tu en as conclu le plus.
C’est d’ailleurs en partie grâce à toi qu’il est toujours gouverneur.
Tu lui donnes régulièrement beaucoup d’argent pour financer sa «~campagne~» (et oui, les dessous de tables sont assez cher ici).
Mais bon, puisqu’il modifie les lois à ton avantage, tu as tout intérêt à rentrer dans son jeu.

Bien entendu certaines personnes parmi la foule commencent à comprendre les magouilles.
Tu as déjà eu pas mal de problème avec la résistance contre le pouvoir sur \textsc{Anacréon}…
C’est un groupe de personne assez peu nombreux, mais totalement incorruptibles et bornés.
Ils vont parfois même jusqu’à sacrifier leur vie dans un idéal absurde ; comme si on pouvait se passer d’un gouverneur tel Paul \textsc{Ledar}…
Jamais l’économie n’en sera plus déséquilibrée.
Si tu voit un moyen simple, et de préférence qui te serait avantageux, pour écraser cette résistance, tu n’hésitera que peu d’instant pour investir une partie de ta fortune dans cette action.

L’avantage à être un proche économique du gouverneur, c’est que l’on voit toujours des personnages important avec qui faire des affaires.
Avec l’arrivée de l’ambassadeur de \textsc{Trantor}, quelques personnes importantes sont ici regroupées.
Si seulement l’un d’eux pouvait être réceptif pour une affaire…

Tu as en effet entendu parler récement d’un traffic d’armes sur la planète proche \textsc{Meneleus}.
Cette affaire t’intéresse grandement : il doit y avoir de grandes quantitées d’argent à gagner par là.

\paragraph{Matériel}
{
Avalga \textsc{Orzon} possède un grand nombre de crédits sur elle.

Il n’y a pas de problème à les utiliser contre toute sorte de choses…
}
}

\player{11}{
\section{Ton personnage : Yovan \textsc{Bunski}}

\paragraph{Description physique}
{
En blouse blanche, grand et mince.
}

Tu es le médecin du vaisseau.
Non seulement tu t’y connais en médecine, mais tu as eu une éducation scientifique assez complète.
Tu as toujours eu un esprit rationnel.

En revanche, le monde dans lequel tu vit (ou plutôt survit en l’occurence) ne te paraît pas rationnel du tout.
D’un côté, il y a ce gouverneur — Paul \textsc{Ledar} — qui règne (il faut le dire, il tiens plus du roi local que du gouverneur) sur \textsc{Anacréon}.
Il est actuellement dans le vaisseau.
Tu as d’ailleurs reçu énormément de menaces sur ta famille : le message est clair, rien ne dois arriver au gouverneur dans ce vaisseau !
De l’autre côté, il y a des gens qui se révoltent contre cet homme et sa politique.
Ils pourraient être des gens très bien, mais ils veulent tout bouleverser par la force…

De ton côté, disons que tu es entre les deux.
Tu n’a pas envie que ta famille soit le moins du monde mêlé là dedans, et tu n’accomplira rien contre le gouverneur si on pourrait remonter à toi en cherchant bien.
D’un autre côté, c’est vrai qu’il serait temps d’en finir avec ce gouvernement de dingue !

Le climat est très tendu — vraiment très tendu.
Tu ne sais pas si tu agira, mais il est probable que tu n’ai pas dit ton dernier mot dans cette histoire.
Le vol spatial est assez long, essayons de voir comment les choses avancent.

Peu après avoir quitté le sol, tu as trouvé dans les toilettes un sachet de poudre brune.
Après analyse, il s’est avéré que cette poudre était de la drogue d’\textsc{Anacréon}.

Cela devient de pire en pire ici…
Bon, il serait de bon ton de trouver qui ici se drogue, ce genre d’informations peuvent toujours servir.
Tu n’aimes pas le chantage, mais il va falloir l’utiliser, tu le crains.
Par contre, il s’agit de le faire discrètement ; sinon, le chantage ne servira à rien de bon.

\paragraph{Capacités spéciales}
{
En tant que médecin, si tu veux faire des analyses sur quelqu’un, quelque chose, etc. tout est possible ; il te suffit d’appeler un maître du jeu.

Tu peux aussi soigner des gens, préparer des somnifères, etc. : de même, il te suffit d’appeler un maître du jeu.
}
}

\player{12}{
\section{Ton personnage : Vanri \textsc{Thrawder}}

Ton personnage vient de la planète \textsc{Smyrno}.
C’est la planète principale d’un secteur voisin, le secteur de \textsc{Smyrno}.

Actuellement, la guerre fait râge sur \textsc{Ifni}, une planète du secteur de \textsc{Smyrno}.
Tu as fait parti des personnes enrôlées plus ou moins de force pour aller servir ton secteur sur cette planète.
Tu y a vu des choses assez horribles, ce qui t’as rendu totalement impassible maintenant.

On t’a vite remarqué là-bas pour tes aptitudes à masquer tes émotions et à agir de façon froide et calculée.
Tu as rapidement grâvi les échelons, en particulier car les missions à risque ne te faisait pas peur.

Au bout de quelques années, le service d’espionnage a fini par s’intéresser à toi.
Tu as tout de suite accepté.
Te voici maintenant espionne de \textsc{Smyrno}.

Le fait que l’Empire commence à s’effacer ici en périphérie a poussé ton gouvernement à s’inquiéter de ses voisins, en particulier le secteur d’\textsc{Anacréon}.
On t’a donc envoyée enquêter sur \textsc{Anacréon} sur les différentes «~magouilles~» politiques qui s’y passe.

Tu as vite réussi à t’intégrer et à voir de nombreuses choses sur la planète.
Tu as tout de suite compris que le gouverneur Paul \textsc{Ledar} était un homme profitant de la corruption pour rester à son poste.
Tu as aussi pris connaissance de la résistance locale contre le gouverneur.
Elle est assez faible, mais composée de quelques personnes capables d’aller jusqu’au bout pour leurs fins.

Tu admires de plus en plus cette \textit{Rébellion}, et tu aimerait te faire connaître d’eux comme alliée.
Malheureusement, cela ne fait pas encore suffisament de temps que tu es ici pour avoir réussi à s’inflitrer.
Mais cela reste dans tes objectifs.

Mais ton but principal reste l’espionnage.
Tu dois ainsi trouver tout ce que ton gouvernement pourrait savoir d’utile sur \textsc{Anacréon}.
Si tu peux, tu dois aussi faire ton possible pour «~couler~» \textsc{Anacréon}, par exemple en déclenchant une guerre civile entre le gouverneur et la résistance.
Bien entendu, tu ne dois en aucun cas de faire connaître comme espionne, ou plus généralement comme originaire de \textsc{Smyrno}.
Tu as tout tes papiers anacreonniens, et tout ce qu’il faut pour te faire passer comme citoyenne d’\textsc{Anacréon}.
Au pire, tu as été entraînée à la torture et aux menaces : on ne pourra rien de faire sortir en te menaçant ; tu préfèrerais mourir que de dévoiler ton origine ou ta véritable identité.

Pour l’instant, tu as réussi à te faire passer comme responsable technique sur ce vaisseau.
Paul \textsc{Ledar} y est présent, tu le sais, et c’est une bonne occasion pour mettre des bâtons dans les roues du gouvernement local.
Personne ici ne semble vraiment s’y connaître en technologie, personne ne te posera de question sur ton rôle donc.
En tant que responsable technique, tu as pu placer quelques mouchards dans le vaisseau.

Tu as également tenté de mettre le bazard dans la sécurité du vaisseau, afin de faciliter tes prochaines actions.
Pour cela, tu as remplacé les pilulles pour la toux de Dacan \textsc{Calrines} (tu avais remarqué qu’il était un peu malade récemment) par des pilules mnémoniques.
Si tout se passe bien, il ne devrait plus se souvenir de quoi que ce soit à son réveil.
Puis tu as laissé quelques «~indices~» dans le vaisseau qui laisserait à penser qu’il organiserait un petit trafic d’armes dans la région de \textsc{Meneleus}, une autre planète du secteur d’\textsc{Anacréon}.
(Il y a maintenant par exemple une caisse pleine d’armes dans sa chambre.)

\paragraph{Matériel}
{
En tant qu’espionne, Vanri \textsc{Thrawder} possède de nombreux petits gadgets d’agent secret.
À chaque utilisation de l’un d’eux, merci de prévenir les maîtres du jeu.

Voici la liste :
\begin{itemize}
	\item Les fameux mouchard dans le vaisseau : tu pourra par exemple demander aux maître du jeu si tu peux réécouter une conversation qui s’est déroulée dans un endroit du vaisseau. (Il est possible qu’il n’y ait pas de mouchard à cet endroit bien entendu…)
	\item Des pilules mnémoniques, somnifères, poisons.
	\item Des armes blanches camouflées. (Attention, effet non garanti : les maîtres du jeu ont tous les pouvoirs)
	\item Des coupes-câbles, et autre matériel permettant de saboter le vaisseau.
	\item Des lunettes infra-rouges.
	\item Un «~tazer~» dissimulé dans un téléphone portable.
\end{itemize}
}
}

\player{13}{
\section{Ton personnage : Dacan \textsc{Calrines}}

\paragraph{Description Physique}
{
Plutôt petit et baraqué.
Ce n’est pas le genre de personne à qui on oserait provoquer.
Tu n’as pas l’air très finaud, mais très menacant.

Habillé en uniforme pseudo-militaire.
}

\paragraph{Le personnage}
{
Le blanc total !

Depuis ce matin, tu ne sais pas ce que tu as, mais tu viens de te relever avec une gueule de bois impossible.
Impossible bien évidement de ce rappeller quoi que ce soit de ce que tu as fait la veille.
Plus embêtant, impossible de se rappeller quoi que ce soit.

Bon. En te réveillant ce matin — un peu perturbé, il faut le dire… — tu finis par faire une petite inspection de tes environs.
Tu t’es réveillé dans une petite chambre très étroite — probablement dans un vaisseau spatial, mais tu ne reconnais pas vraiment le vaisseau (de toute façon, ils se ressemblent tous…).
Un uniforme est jeté sur la chaise près de ton lit.
Il te va parfaitement ; tu le met.
La plaque inscrite dessus porte le nom de Dacan \textsc{Calrines}… Cela ne te dit rien, mais tu ne te souviens même plus de ton propre nom.
Bon, disons que c’est le tiens.

Tu te lève, et remarque que ton oreiller a une texture étrange.
Ça ne ressemble pas aux oreillers habituels dans ce genre de vaisseau.
Tu l’éventre et trouve quelques petits sac… À peine 10cm par 5cm, transparents, tous remplis d’une poudre brune.
Tu ne vois pas vraiment ce que ça pourrait être, mais tu ne préfère pas trop savoir.
De la drogue ? Des nanites inactivées ? Tu ne vois pas.

Bon, continuons l’inventaire… Sous le lit ?
Il se trouve quelques affaires personnelles, ainsi qu’une caisse bien bazardée au fond.
Lorsque tu la sort de sous le lit, cela laisse apparaître une grande quantité de crédit.
Ça, par contre, tu te souviens que tout cela représente beaucoup d’argent pour quelqu’un n’ayant qu’un uniforme comme le tiens.
Tu ouvres non sans difficulté la caisse : elle est remplie de blaster !

Ça commence à devenir vraiment louche, là.
Bon, tu préfère ne pas en savoir plus pour l’instant, essayons plutôt d’explorer le reste du vaisseau.

Sur ta porte, un post-it te rappelle d’aller voir \textsc{Idaski} pour un briefing (enfin, tu imagines que c’est ça, c’est de la prise de note assez elliptique).
Une heure est indiquée, tu regardes ta montre : c’est dans à peine 20 minutes !
Si on ne veux pas trop te poser de questions, et pour en savoir facilement plus, il vaudrait mieux que tu trouve cette \textsc{Idaski} assez vite !
}

\paragraph{Matériel}
{
Ton personnage possède sur lui un blaster (les tirs sont léthaux sur un humain si envoyé aux bons endroits) et un communicateur.

À chaque utilisation d’un de ces éléments, merci de prévenir un maître du jeu.
}
}


\end{document}
